\section{Infraestructura y la carrera del cómputo}

\subsection{Clústeres a hiperescala}

La infraestructura de IA está atravesando una velocidad de expansión sin precedentes en la historia humana. Para poner esto en perspectiva: GPT-4 usó 20,000 GPU A100 ($\sim$15-20 MW), equivalente a un centro de datos estándar. Apenas dos años después, la escala de un solo clúster ya se ha multiplicado por 10.

\begin{table}[H]
\centering
\begin{tabular}{lrl}
\toprule
\textbf{Organización} & \textbf{Cantidad de GPU} & \textbf{Notas} \\
\midrule
xAI (Memphis) & 200K H100 + 100K H20 & Clúster individual más grande; meta de Elon: 1M \\
Meta & $\sim$128,000 & 400K+ en total \\
OpenAI & $\sim$100,000 & \\
Google (TPU) & Total más grande & Fibra óptica entre Iowa, Nebraska y Ohio \\
Anthropic + Amazon & 400K Trainium 2 & En construcción \\
\bottomrule
\end{tabular}
\caption{Escala de los clústeres globales de cómputo de IA (principios de 2025)}
\end{table}

El consumo eléctrico de los centros de datos pasará del 2-3\% de la electricidad de Estados Unidos a más del 10\% hacia 2028-2030. La próxima generación de GPU Blackwell alcanza un consumo de 1200W (frente a los 700W actuales de Hopper); la meta para el próximo año es un clúster de 500,000-700,000 GPU. La meta de Elon es un millón de GPU: \textit{«Nunca dudo de Elon.»} (Dylan)

\begin{knowledgebox}{El problema de los picos de potencia en el entrenamiento}
El entrenamiento genera una demanda eléctrica en forma de picos: potencia alta cuando las GPU calculan $\to$ las GPU quedan inactivas y con baja potencia mientras se intercambian los pesos. Esta carga pulsante resulta muy poco amigable para la red eléctrica.

Meta liberó por accidente un operador de PyTorch: \texttt{PowerPlant.no\_blowup = 1}, que hace que las GPU calculen números sin sentido durante los periodos de inactividad para suavizar la curva de potencia. \textit{«Es muy fácil terminar quemando el equipo.»} (Dylan)

La solución de Elon es más elegante: los paquetes de baterías Tesla Mega Pack absorben las transiciones bruscas de potencia.
\end{knowledgebox}

\subsubsection{La revolución del enfriamiento líquido}

El enfriamiento por aire solía ser la solución estándar, pero la densidad térmica de 1200W de Blackwell lo vuelve inviable. Google TPU ha usado enfriamiento por agua desde hace años. Elon fue pionero en el enfriamiento líquido de GPU a gran escala en Memphis: 90 unidades de enfriamiento por agua del tamaño de contenedores rodean el perímetro de las instalaciones. La próxima generación de NVIDIA exigirá enfriamiento líquido.

Un beneficio colateral del enfriamiento líquido: los chips pueden colocarse más juntos $\to$ distancias físicas más cortas $\to$ interconexiones de mayor velocidad. Esto resulta fundamental para la eficiencia del entrenamiento.

\subsection{Detalles del proyecto Stargate}

El proyecto de OpenAI y Oracle en Abilene, Texas: potencia a plena carga de 2.2 GW, superior al consumo eléctrico de la mayoría de las ciudades.

\begin{importantbox}{La escala y el financiamiento de Stargate}
\begin{itemize}
    \item El costo total de propiedad de la Fase 1 es de aproximadamente \$100B (\$50B de inversión real + \$50B de costos operativos)
    \item Primera etapa: $\sim$\$5-6B en servidores + \$1B en el centro de datos (Oracle ya está construyendo)
    \item Fuentes de financiamiento: SoftBank (posiblemente \$25B), Oracle, MGX (Emiratos Árabes Unidos, fondo de IA de \$1.5T pero «con la tinta aún sin secar»)
    \item El propio compromiso de OpenAI es de \$19B, pero solo cuenta con \$6B en efectivo + \$4B en deuda
\end{itemize}

Elon tiene razón: \textit{«El dinero no existe.»} Pero Dylan cree que el dinero terminará por llegar. El papel de Trump es la desregulación (centros de datos en terrenos federales, simplificación de permisos), no aportar financiamiento. \textit{«Trump es un hype man... reducir la regulación les permite construir más rápido.»} (Dylan)
\end{importantbox}

\begin{figure}[H]
\centering
\includegraphics[width=0.92\textwidth]{frame-04-stargate.jpg}
\caption{Discusión sobre el gasto de capital y la escala de infraestructura de Stargate\protect\footnotemark}
\end{figure}
\footnotetext{Intervalo de tiempo del video: 03:43:00--03:43:12.}

\subsection{Electricidad, enfriamiento y restricciones físicas}

El entrenamiento genera una demanda eléctrica en forma de picos: potencia alta cuando las GPU calculan, baja potencia y GPU inactivas mientras se intercambian los pesos. Meta liberó por accidente un operador de PyTorch: \texttt{PowerPlant.no\_blowup = 1}, que hace que las GPU calculen números falsos durante los periodos de inactividad para suavizar la potencia.

\textit{«Es muy fácil terminar quemando el equipo.»} (Dylan)

La solución de Elon: los paquetes de baterías Tesla Mega Pack absorben las transiciones bruscas de potencia. La revolución del enfriamiento líquido: Google TPU ya usa enfriamiento por agua, y Elon fue pionero en el enfriamiento líquido de GPU a gran escala en Memphis (90 unidades de enfriamiento por agua del tamaño de contenedores). La próxima generación NVIDIA Blackwell exigirá enfriamiento líquido. El enfriamiento líquido permite colocar los chips más juntos, logrando interconexiones de mayor velocidad.

\subsection{Contrabando de GPU y reglas de difusión}

\begin{warningbox}{El contrabando de GPU ya es una realidad}
ByteDance es el «contrabandista» más grande: renta más de 500,000 GPU en todo el mundo a través de Oracle, Google y diversos proveedores pequeños de nube. El contrabando físico también existe: hay quienes viajan en primera clase desde el aeropuerto de San Francisco a Shanghái con cajas de GPU de SuperMicro.

Se estima que en 2024 entre 200 y 300 mil GPU llegaron a China a través de empresas pequeñas de Singapur y Malasia. Las nuevas reglas de difusión de IA limitan a las entidades vinculadas con China a rentar menos de 2000 GPU o comprar menos de 1500 GPU.

\textit{«Las GPU superarán a las drogas y las armas como el producto de contrabando de mayor valor por kilogramo.»} (Dylan)
\end{warningbox}

\subsection{Google TPU y el monopolio de NVIDIA}

El TPU de Google es excelente para uso interno (optimizado para Búsqueda, YouTube y Gemini), pero su pila de software no se ha publicado de forma abierta. Se trata de un problema organizacional y no técnico: Google Cloud, el equipo de TPU, DeepMind y el equipo de Búsqueda operan de manera independiente entre sí, sin una estrategia externa unificada. Toda la cultura de NVIDIA está orientada a servir a clientes externos; Google sirve a clientes internos.

Un ejemplo concreto revela las desventajas de esta optimización interna: el tamaño del vocabulario del modelo Gemma está optimizado para TPU (que cuentan con unidades grandes de multiplicación de matrices), pero resulta muy poco eficiente al ejecutarse en GPU. Solo cuando los investigadores de Google salen a fundar empresas descubren lo difícil que es la infraestructura externa: \textit{«Los investigadores dejan Google, fundan una empresa, descubren que la infraestructura es difícil, y luego regresan.»} (Dylan)

El caso de AMD: el hardware es aceptable, pero el software (ROCm) es pésimo hasta lo indignante. \textit{«Somos los que más bugs reportamos a AMD: ¿por qué somos nosotros, SemiAnalysis, quienes reportamos más bugs?»} (Nathan) Que una empresa de análisis de semiconductores esté reportando bugs de software a un fabricante de chips muestra lo endeble que es el ecosistema de software de AMD.

La situación de Intel es aún más difícil: perdió el liderazgo en el proceso de fabricación, no ha logrado ningún éxito en chips de IA, su director ejecutivo fue destituido y sigue perdiendo participación de mercado en PC y servidores. El rendimiento de fabricación de Samsung también está en declive.

El escenario más extremo: el mundo podría terminar dependiendo únicamente de TSMC para la investigación y el desarrollo de chips de vanguardia. Esto significa que una isla, situada a unos cientos de kilómetros de la costa de China continental, carga con el destino de la tecnología global.

\subsection{Resumen de la sección}

La infraestructura de IA está atravesando una expansión sin precedentes: proyectos del calibre de Stargate representan inversiones del orden de \$100B. La electricidad, el enfriamiento y el contrabando de GPU son las nuevas restricciones físicas. El monopolio de NVIDIA no podrá revertirse en el corto plazo: el software de AMD es demasiado deficiente, Intel está en declive y el TPU de Google no está disponible externamente.
